\documentclass[11pt]{article}
\usepackage[margin=1in]{geometry}
\usepackage{amsmath,amssymb}
\usepackage{booktabs}
\newcommand{\Hd}{\widetilde{H}_d}
\newcommand{\Ht}{\widetilde{H}}
\begin{document}
\begin{center}
\textbf{Distributed GLR: factorization, spectral operations, and memory layouts}\\[2pt]
\small Design record for the GLR engine of \texttt{lgpsf-hessian} (the
\texttt{lgh\_glr\_*} API; ScaLAPACK backend).  Originally written 2026-08-16 as the
design record of the engine's first production deployment, a large-scale ice-sheet
inversion; the outer factors $Z, M$ are the library's prior interface
($R = Z M^{-1} Z$, \texttt{lgh\_prior\_t}).  Supersedes the replicated-$T$ structure
documented in \texttt{glr-woodbury-formulas.pdf} (which the library keeps as its
replicated backend).
\end{center}

\section{Operators and factorization}

$N$ dofs, distributed by the application's spatially-local mesh partition over $P$
ranks.  $M$ = diagonal lumped mass.  $Z$ = the prior's square-root factor (e.g.\ a
shifted Laplacian at a reference regularization scale); $Z^{-1}$ available as a
linear operator (multigrid-preconditioned Chebyshev, \texttt{lgh\_prior\_create\_mat}).
Prior precision $H_r = Z M^{-1} Z$ (bilaplacian-like).  $\Hd$ = symmetrized lgpsf
approximation of the misfit Hessian (stage 1, \texttt{lgh\_fit\_*}), available
through matvecs on mesh-partitioned vectors.  $c$ = regularization /
Levenberg--Marquardt weight.
\begin{align}
\Ht(c) &= \Hd + c\,H_r
        = Z M^{-1/2}\left( F + cI \right) M^{-1/2} Z,\\
F &= M^{1/2} Z^{-1} \Hd\, Z^{-1} M^{1/2} = F^T
        \qquad (\text{whitened prior-preconditioned misfit Hessian; never formed}).
\end{align}

\emph{Keep $c$ and the regularization scale $\alpha$ explicit; do not absorb them
into $Z$.}  Reasons: (i) one eigendecomposition of $F$ serves every $c$ through the
spectral filters below (the library's operations all take $c$ per call), whereas $c$
absorbed into $Z$ would force a rebuild per shift; (ii) for the bilaplacian with
$Z = \alpha A$, changing $\alpha$ rescales eigenvalues exactly with unchanged
eigenvectors, $F(\alpha') = (\alpha/\alpha')^2 F(\alpha)$ --- a free update, visible
only if $\alpha$ is a scalar in the code.

\section{Target operations}

\begin{enumerate}
\item Remove spurious negative eigenvalues of $F$ (approximation artifacts):
      FLIP treatment $\lambda_i \leftarrow |\lambda_i|$
      (\texttt{LGH\_TREAT\_FLIP}, the default).  Do not discard large-magnitude
      negatives --- they are genuine flipped-curvature modes, and clipping them to
      zero leaves $\mathcal{O}(1+|\lambda|)$ outliers in the preconditioned spectrum.
\item Solves with the (positively modified) $\Ht$: the deterministic MAP phase
      (\texttt{lgh\_glr\_solve}).
\item Square-root and inverse-square-root \emph{factor} actions for MCMC sampling
      (\texttt{lgh\_glr\_factor}, \texttt{lgh\_glr\_sample}).
\item Log-determinant differences $\log\det\Ht_1 - \log\det\Ht_2$ for detailed
      balance with location-dependent proposal precision (\texttt{lgh\_glr\_logdet}).
\end{enumerate}

\section{Randomized eigendecomposition of $F$}

Sketch width $\ell$ (target rank $r \le \ell$ after truncation; continental-scale
expectation in the originating application $r \sim 10^4$--$2\times 10^4$, must scale
to $5\times 10^4$+).
\begin{align}
\Omega_{ij} &= \mathrm{hashnormal}(\mathrm{seed}, i, j)
   \quad (\text{drawn locally from global index; partition-independent})\\
Y &\leftarrow F\,\Omega; \qquad
   \text{for } p = 1,\dots,q: \ \ Q \leftarrow \mathrm{orth}(Y),\ Y \leftarrow FQ\\
Q &\leftarrow \mathrm{orth}(Y) \qquad
   (\text{panel-blocked BCGS2; no } \ell\times\ell \text{ Gram is ever materialized})\\
T &= \mathrm{sym}(Q^T F Q) \qquad
   (\text{formed in column panels, reduced directly into 2D block-cyclic layout})\\
T &= V \Lambda_{\mathrm{raw}} V^T \qquad
   (\text{distributed dense eigensolver: ScaLAPACK \texttt{pdsyevd}, ELPA later;}\\
&\qquad\quad V \text{ \emph{stays in place}, block-cyclic; } \Lambda \text{ replicated})\notag\\
U &= Q V_{:,\mathcal{K}} \quad \text{\emph{implicit}: } Uw = Q(Vw),\ \ U^T x = V^T(Q^T x);
   \qquad \lambda_i = |\lambda_{\mathrm{raw},i}| > \tau \ \ (\text{FLIP + cut})
\end{align}
$U^T U = I_r$; $F \approx U \Lambda U^T$.  The coefficient traffic per $U$-apply is
$\mathcal{O}(\ell)$ numbers --- negligible.  At large $\ell$ the implicit form is
forced, not merely elegant: materializing $U = QV$ with block-cyclic $V$ would
require either replicating $V$ ($8\ell^2$ bytes per rank; 20\,GB at
$\ell = 5\times 10^4$) or a heavyweight mixed-layout distributed GEMM.

\emph{Incremental rank growth (first-class operation: \texttt{lgh\_glr\_extend}).}
Since $r$ is not known a priori: extend by a block of $k_{\mathrm{new}}$ fresh
hashed columns (global column index continues past the current $\ell$, so the
extension is deterministic and partition-independent), apply $F$,
BCGS2-orthogonalize against the existing $Q$, border $T$ with
$Q^T (F Q_{\mathrm{new}})$, re-eigensolve.  Stop when the sketch's smallest
$|\lambda|$ (reported as \texttt{next\_abs}) falls below the truncation cut $\tau$
with margin --- the unknown rank becomes a stopping criterion instead of a gamble on
$\ell$.  (Evidence this matters: in the originating application, a Pine Island
Glacier sketch at $\ell = 2200$ was rank-limited --- 2136 of 2200 modes survived
the cut.)

\section{Spectral operations}

With the truncated tail of $F$ treated as $0$ (the GLR model itself), any spectral
function acts as
\begin{equation}
f(F + cI) \;=\; U\!\left[ f(\Lambda + cI) - f(c)\,I_r \right]\! U^T \;+\; f(c)\, I
\qquad (\texttt{lgh\_glr\_filter}).
\label{eq:specfun}
\end{equation}
Only the actions of $U$ and $U^T$ are required.

\paragraph{Solve (MAP phase).}  $f(x) = 1/x$ in \eqref{eq:specfun} gives
\begin{equation}
\Ht(c)^{-1} = Z^{-1} M^{1/2}
  \tfrac{1}{c}\!\left( I - U D_c U^T \right) M^{1/2} Z^{-1},
\qquad D_c = \mathrm{diag}\!\left( \frac{\lambda_i}{\lambda_i + c} \right),
\end{equation}
(\texttt{lgh\_glr\_solve}), which at $c = 1$ is exactly the Woodbury apply of the
engine's originating production code (continuity check).  The inverse is the one
case where the whitened spectral function composes exactly with the outer factors,
because it is the inverse of a product.

\paragraph{Square-root factors (MCMC).}  The sandwich is \emph{not} a similarity
transform, so the symmetric square root $\Ht^{1/2}$ is not cheaply available --- and
is not needed.  Sampling requires a factor:
\begin{equation}
G = Z M^{-1/2} (F + cI)^{1/2}, \qquad \Ht(c) = G\,G^T,
\end{equation}
with $(F+cI)^{\pm 1/2}$ from \eqref{eq:specfun}.  All four factor actions
($G,\ G^T,\ G^{-1},\ G^{-T}$; \texttt{lgh\_glr\_factor}) reduce to $U$/$U^T$
applies, $M^{\pm 1/2}$ scalings, and $Z$ matvecs or solves.  A draw from the
proposal $\mathcal{N}(m, \Ht^{-1})$ is $x = m + G^{-T}\xi$,
$\ G^{-T} = Z^{-1} M^{1/2} (F + cI)^{-1/2}$, $\ \xi \sim \mathcal{N}(0, I)$
(\texttt{lgh\_glr\_sample}).

\paragraph{Log-determinant differences.}  $Z$ and $M$ are state-independent and
cancel between two builds:
\begin{equation}
\log\det \Ht_1(c) - \log\det \Ht_2(c)
  = \sum_k \log\!\big(\lambda^{(1)}_k + c\big) - \sum_k \log\!\big(\lambda^{(2)}_k + c\big)
\end{equation}
where the per-build quantity to compute (\texttt{lgh\_glr\_logdet}) is
$\sum_k [\log(\lambda_k + c) - \log c]$ over the \emph{kept, treated} eigenvalues:
truncated modes contribute $0$, so the sum is finite and truncation-consistent, and
the common $Nc$-dependent constant cancels in the difference.
\emph{Detailed-balance consistency:} Metropolis--Hastings needs the log-determinant
of the proposal precision \emph{actually sampled from}.  Sampling and
log-determinant both read the same treated $\Lambda$, so truncation and FLIP do not
threaten correctness --- they only affect proposal quality (acceptance rate).

\section{Shapes and memory layouts}

Feature-space principle: the reduced space $\mathbb{R}^\ell$ is distinct from the
function space $\mathbb{R}^N_M$ and its dual, and owns its own partition --- by
feature index (2D block-cyclic), never by mesh geometry.  Small feature-space
objects ($\Lambda$, filter diagonals) are replicated; that is a storage convention,
not a leak between spaces.

\begin{center}
\begin{tabular}{llll}
\toprule
Symbol & Shape & Space & Layout \\
\midrule
$v, x, w_{\mathrm{full}}$ & $N$ & function / dual & mesh partition ($N_p$ per rank) \\
$M^{\pm 1/2}$ & $N$ (diag) & --- & mesh partition \\
$Z$ & $N \times N$ sparse & dual $\leftarrow$ function & mesh partition; $Z^{-1}$ via MG-Chebyshev \\
$\Hd$ & $N \times N$ sparse & dual $\leftarrow$ function & mesh partition (MPIAIJ) \\
$\Omega, Y, Q$ & $N \times \ell$ dense & whitened & mesh-partitioned row slabs ($N_p \times \ell$) \\
BCGS2 coefficients & $\ell \times b$ & --- & transient, replicated per panel (discarded) \\
$T$ & $\ell \times \ell$ & feature & 2D block-cyclic on $p_r \times p_c$ grid; never replicated \\
$V$ & $\ell \times \ell$ & feature & block-cyclic, \emph{stays in place} after eig \\
$\Lambda$, $D_c$, filters & $r$ (diag) & feature & replicated (small) \\
$U = Q V_{:,\mathcal{K}}$ & $N \times r$ & whitened $\leftarrow$ feature & \emph{implicit composition; never formed} \\
$w$ (reduced vectors) & $\ell$ or $r$ & feature & conformal with $T$'s layout; \\
 & & & gathered ($\mathcal{O}(\ell)$ traffic) inside $U$-applies \\
$c, \alpha, \tau$ & scalars & --- & replicated \\
\bottomrule
\end{tabular}
\end{center}

Memory at the ambition point ($N = 10^7$, $\ell = 5\times 10^4$, $P = 100$ nodes):
$T, V$ are $20\,$GB \emph{aggregate} ($\sim$0.2\,GB/node, block-cyclic) --- no longer
a constraint.  The tall blocks dominate: $Q$ is $4\,$TB aggregate ($40\,$GB/node),
and the sweep buffers must be managed (overwrite in place, panel the $F$ sweeps).
At $N = 10^6$ everything is comfortable.

\section{Design rationale (brief)}

\begin{itemize}
\item \textbf{All three $\ell\times\ell$ objects go distributed, not just $T$.}
  The orthogonalization Gram matrix has the same $\mathcal{O}(\ell^2)$ replicated
  footprint as $T$; panel-blocked BCGS2 avoids materializing it at all, is more
  robust than CholQR for power-iterated (ill-conditioned) blocks, and makes
  incremental column-block extension natural.
\item \textbf{$T$ is formed directly into its block-cyclic home}, in column panels
  (transient $\ell \times b$ per rank), never assembled replicated.
\item \textbf{Implicit $U$} avoids both the $N\ell r$ materialization GEMM and any
  replication of $V$; per-apply cost is essentially unchanged (GEMV over $\ell$
  columns of $Q$ vs.\ $r$ columns of $U$, plus $\mathcal{O}(\ell)$ coefficient
  traffic).
\item \textbf{Incremental growth} turns unknown rank into a stopping criterion;
  supporting $r \gtrsim 5\times 10^4$ cheaply is a core advantage of lgpsf-based
  approximation over conventional low-rank methods and must not be blocked by a
  fixed-$\ell$ pipeline.
\item \textbf{Eigensolver}: ScaLAPACK \texttt{pdsyevd} first (ships with MKL on
  large clusters; reference build available everywhere via PETSc), ELPA/SLATE
  behind the same seam when $\ell \gtrsim 3\times 10^4$ makes it pay.  Eigenvector
  signs / near-degenerate rotations depend on the process grid, so cross-$P$ (and
  cross-backend) comparisons use invariants only ($\Lambda$, operator applies,
  logdet).
\item \textbf{Determinism}: hashed $\Omega$ (and hashed extension columns) keep the
  sketch partition-independent; a dedicated study of $T$ on the originating problem
  showed it is a generic dense symmetric matrix (no banding, no hierarchical
  structure, full-rank off-diagonal blocks), so a dense distributed eigensolver is
  the right primitive --- there is nothing to exploit by compression.
\end{itemize}
\end{document}
